% !TEX root = userguide.tex

\def\headdate{14th October 2021}

\section{Persistence-of-straining}

(This section and the examples have been contributed by Jang Min Park of the Yeungnam University.)\medskip

In this section the persistence-of-straining concept
developed by Thompson and Souza Mendes \cite{Thompson2005}
is introduced briefly,
and some examples of its computation are presented.

\subsection{Basics}
\label{sec:persist_basic}
The rate-of-strain and the rate-of-rotation tensors are
defined as 
$\ten{D} = \dfrac{1}{2} \left( \ten{L} + \ten{L}^T \right)$ and
$\ten{W} = \dfrac{1}{2} \left( \ten{L} - \ten{L}^T \right)$, respectively,
where
$\ten{L}= \left( \nabla \vec{u} \right)^T$ is the velocity gradient tensor.

Since $\ten{D}$ is symmetric, 
its eigenvalues, namely $\lambda^{(p)}, \, p=1, 2, 3,$ are real
and an orthonormal set of eigenvectors 
$\vec{e}^{(p)}, \, p = 1, 2, 3,$ can be obtained such that
\begin{equation}
\ten{D} \cdot \vec{e}^{(p)} =
\lambda^{(p)} \vec{e}^{(p)}
\quad \,
\left( p = 1, 2, 3 \right)
\end{equation}

A tensor $\ten{\Omega}$ representing the rate-of-rotation of 
the principal axis (eigenvectors) of $\ten{D}$,
is defined in such a way that
\begin{equation}
\dot{\vec{e}}^{(p)}
=
\ten{\Omega}
\cdot
\vec{e}^{(p)}
\end{equation}
where overdot ($\dot{\quad}$) represents the material time derivative of its argument.
%
Since $\vec{e}^{(p)} \cdot \vec{e}^{(q)} = \delta_{pq}$, 
one can find that $\ten{\Omega}$ is antisymmetric.
%
Let us write the tensor $\ten{\Omega}$ in the basis defined by eigenvectors of $\ten{D}$ as
\begin{equation}
\ten{\Omega} =
\sum_{p=1}^3 \sum_{q=1}^3
{\Omega}_{pq}^{*}
\vec{e}^{(p)} \otimes \vec{e}^{(q)}
\end{equation}
In the followings, asterisk symbol ($^*$) in the matrix representation means
the basis is the eigenvectors of $\ten{D}$,
while no symbol means the basis is of a fixed reference coordinate system.
%
An explicit expression for ${\Omega}_{pq}^{*}$ can be
written as
\begin{equation}
\begin{split}
{\Omega}_{pq}^{*}
&=
\vec{e}^{(p)}
\cdot
\ten{\Omega}
\cdot
\vec{e}^{(q)}
=
-
\frac
{\vec{e}^{(p)} \cdot \dot{\ten{D}} \cdot \vec{e}^{(q)}}
{\lambda^{(p)} - \lambda^{(q)}}
=
-
\frac
{\dot{D}_{pq}^{*}} 
{\lambda^{(p)} - \lambda^{(q)}}
, \quad 
(p \ne q)
\\
{\Omega}_{pq}^{*}
&= 0, \quad
(p = q)
\end{split}
\end{equation}
It should be noted that the above equation assumes 
$\lambda^{(p)} \ne \lambda^{(q)}$.
%
In the degenerate case of 
$\lambda^{(p)} = \lambda^{(q)}$ which are different from the third eigenvalue,
${\Omega}_{pq}^{*}$ is not uniquely defined.

The persistence-of-straining tensor $\ten{P}$ is defined as
\begin{equation}
\ten{P} =
\ten{D} \cdot \ten{\overline{W}} 
-
\ten{\overline{W}} \cdot \ten{D}  
\end{equation}
where
$\ten{\overline{W}} = \ten{W} - \ten{\Omega}$ represents
the relative rate of rotation with respect to the principal axis
of $\ten{D}$.
%
In the basis of eigenvectors of $\ten{D}$,
the matrix of $\ten{P}$
can be written as
\begin{equation}
\begin{split}
\mat{\ten{P}}^{*} 
& =
\left[
\begin{matrix}
0 & 
\left(\lambda^{(1)}-\lambda^{(2)}\right) W_{12}^* & 
\left(\lambda^{(1)}-\lambda^{(3)}\right) W_{13}^* \\
\left(\lambda^{(1)}-\lambda^{(2)}\right) W_{12}^* & 
0 &
\left(\lambda^{(2)}-\lambda^{(3)}\right) W_{23}^* \\
\left(\lambda^{(1)}-\lambda^{(3)}\right) W_{13}^* & 
\left(\lambda^{(2)}-\lambda^{(3)}\right) W_{23}^* &
0
\end{matrix}
\right]
\\
& +
\left[
\begin{matrix}
0 & 
\dot{D}_{12}^* & 
\dot{D}_{13}^* \\
\dot{D}_{12}^* & 
0 &
\dot{D}_{23}^* \\
\dot{D}_{13}^* & 
\dot{D}_{23}^* &
0
\end{matrix}
\right]
\end{split}
\end{equation}
Therefore,
$\ten{P}$ is well defined even if two eigenvalues are equal so $\ten{\Omega}$
is not uniquely defined.
It might be mentioned that diagonal part of matrix $\mat{\ten{D}}^*$ is not involved in 
matrix $\mat{\ten{P}}^*$.

The flow classification parameter $R$ based on 
the persistence-of-straining is defined as
\begin{equation}
R = \dfrac{\sqrt{\dfrac{1}{2} \tr {\ten{P}^2}} }{\tr {\ten{D}^2}}
\end{equation}
This parameter represents the strength of \emph{non}-persistence of the straining such that
$R = 0$ corresponds to the complete persistence of straining (extensional flows), 
while $R \to \infty $ corresponds to the minimum persistence of straining (nearly rigid body motion).
For a viscometric flow such as a simple shear flow, $R = 1$.

Some more details regarding the above equations are available in \cite{Park2020}

\subsection{Shear and shearfree flows in a cubic domain}
This program does not solve any problem but it only computes
the flow classification parameter for a given velocity field on a cubic domain.
Either shear or shearfree flow can be applied depending on the
choice of the function number defined in the program.
For shear flow, the flow classification parameter should be unity,
while it should be zero for extensional flow, as mentioned in \ref{sec:persist_basic}.

The example file can be obtained by typing:
\begin{quote}
  \texttt{getexample persist\_strain}
\end{quote}
and you can run the example by typing:
\begin{quote}
  \texttt{persist1}
\end{quote}

\subsection{2D Taylor-Couette flow}
\label{sec:persist_taylorcouette}
This program first solves the Stokes problem in a segment of a Taylor-Couette flow,
and then it computes the flow classification parameter.
When the angular velocities of inner and outer cylinders are set to be different,
the flow is a simple viscometric flow and therefore the flow classification parameter
should be unity.
On the other hand, if the inner and outer cylinders have the same angular velocity, 
the fluid motion is obviously a rigid body motion, 
and the flow classification parameter should approach infinity.

The example file can be obtained by typing at the command line:
\begin{quote}
  \texttt{getexample taylor\_couette}
\end{quote}
and you can run the example by typing:
\begin{quote}
  \texttt{taylor\_couette2 < input\_taylor\_couette2.txt}
\end{quote}
Note, Gmsh needs to be installed on the system
since the program executes gmsh from the command line
to generate a mesh file (\texttt{mesh2.msh})
using the geometry file (\texttt{mesh2.geo}).
The program will create a directory \texttt{taylor\_couette2\_files}
where the result file will be written.

